\documentclass[../main.tex]{subfiles}
\begin{document}

\subsection{Background}

With the expansion of the unmanned vehicle (land, air, surface and underwater)
in many battlefield and disaster scenarios, there is a demand for methods of
controlling these. Whilst autonomy is garnering a lot of attention in this way,
there is a need for the expertise and skill that a human controller brings,
whilst standing them off from the dangers of hazardous environments.

In response to this, development of a remote control system is advisable
.There have been numerous solutions to this in the past, but often they are static and
they require infrastructure.

\subsection{Problem Statement}

Autonomy in workboats is as such becoming increasingly mature and is being deployed to many different platforms.
Until a level of assurance in the autonomous controller's capabilies and reliability is reached (if such a level will ever be reached), human oversight of the vessel will be required in some from.
Even in situations where autonomy is not the end goal for the vessel, there is a need for a method of removing the crew from the vessel in hazardous conditions.
This can reduce the chance and/or severity of accidents (man overboard, injuries from machinery, loss of platform) whilst potentially improving quality of life for those crews
(e.g.\ protection from the elements).
Moreover, there is an opportunity for workboats to be constructed without a wheelhouse and on board controls, which could reduce the cost, complexity, and displacement of the platform.
In all of these situations remote control of the vessel is desirable: to keep watch over autonomy or remove the crew from the vessel.
As mentioned above, Workboat Code edition 3 Annexe 2 \cite{ShipStandards2025} prescribes requirements for an ROUV, but does not prescribe solutions to this issue.

In designing a remote control unit (RCU), the manner in which the operator issues commands and receives information is important.
The RCU will have to be provide situational awareness that is at least as good as a crew member would have aboard the vessel.
It also must be able to provide the same functionality and level of control as those controls normally found aboard the platform.
Additionally, it must be capable of connecting to the vessel in a reliable and secure way.
Such a communication channel must be capable of providing an appropriate SIL, and must permit the transmission of commands, information, and alerts simultaneously.
There a number of approaches to this that are dependant on the setting in which the controller would be used.

There are also specific needs that an RCU will need to meet introduced by the remote control.
The platform can only accept one controller at a time to avoid complexity in controller hierarchy.
As such, the RCU will need to manage assuming control of the vessel, both from a ``cold start'' and from an exant controller.
Additionally, the RCU will need to manage different methods of control behaviour (open water operation, confined spaces, and dynamic station keeping) whilst also being able to manage a degraded
communication environment (such as in inclement weather).

Managing ROUV subsystems is also core functionality that must be met.
Control of the propulsion and power systems is necessary for controlling movement and for enabling a ``cold start'' scenario.
COLREGs 1972 mandates shapes, lights, and sounds that must be managed depending on the boat's situation and movements, and must be controlled as well.
Additionally, the emergency system aboard the vessel must be controlled as well.
This includes both alerts to the operator and the emergency stop functionalities.
To that end, the vessel must also comply with all relevent legislation, such as that pertaining to the ROUV and also to those regarding radio communications.

\subsection{System Narrative}

\subsection{Aims and Objectives}
